\begin{figure*}[!tbp]
\centering
\resizebox{\linewidth}{!}{%
\begin{tikzpicture}[
    font=\sffamily, >=Latex,
    input/.style={draw=ngBlue, top color=ngBlue!8, bottom color=ngBlue!28,
                  rounded corners=4pt, align=center, minimum width=34mm,
                  minimum height=12mm, font=\small},
    interp/.style={draw=ngPurple, top color=ngPurple!10, bottom color=ngPurple!30,
                  rounded corners=4pt, align=center, minimum width=52mm,
                  minimum height=20mm, font=\small,
                  blur shadow={shadow blur steps=5}},
    outcome/.style={draw=ngGreen, top color=ngGreen!10, bottom color=ngGreen!30,
                  rounded corners=4pt, align=center, minimum width=36mm,
                  minimum height=12mm, font=\small},
    note/.style={draw=ngOrange, top color=ngOrange!10, bottom color=ngOrange!30,
                  rounded corners=3pt, align=center, minimum width=42mm,
                  minimum height=14mm, font=\footnotesize},
    design/.style={draw=ngRed, top color=ngRed!10, bottom color=ngRed!30,
                  rounded corners=3pt, align=center, minimum width=44mm,
                  minimum height=14mm, font=\small},
    arr/.style={-{Latex[length=2.8mm]}, line width=0.9pt, draw=ngSlate},
    fbarr/.style={-{Latex[length=2.6mm]}, line width=0.9pt,
                  draw=ngRed!75, dashed}
]
% --- Top row: inputs (left), interpretation (centre), outcomes (right).
\node[input]   (cat)    at (0   , 3.6)  {Data category\\\emph{personalness}};
\node[input]   (not)    at (0   , 1.8)  {Disclosure notice\\\emph{mention}};
\node[input]   (bnd)    at (0   , 0.0)  {Policy boundary\\\emph{coll. / shar.}};

\node[interp]  (interp) at (6.0 , 1.8)  {%
    \textbf{Contextual interpretation}\\[2pt]
    \scriptsize app function $\cdot$ benefit\\
    \scriptsize recipient $\cdot$ control $\cdot$ alternative};

\node[outcome] (nec)    at (11.5, 2.7)  {Necessity\\judgment};
\node[outcome] (dec)    at (11.5, 0.9)  {Privacy\\decision};

% --- Bottom row: interpretive burden and design implications.  The two
%     boxes are made narrower so a visible arrow can run between them,
%     and the row is placed at y=-1.6 (closer to the top row) for
%     visually balanced spacing.
\node[note]    (burden) at (4.3 ,-1.6)  {Interpretive burden\\(ambiguous categories\\\& boundaries)};
\node[design]  (design) at (11.5,-1.6)
    {\textbf{Design implications}\\[2pt]
     \scriptsize Why $\cdot$ Where $\cdot$ Who\\
     \scriptsize $\cdot$ Control $\cdot$ Flag};

% --- Forward arrows (inputs -> interp -> nec -> dec).  TikZ resolves
%     the endpoints against each box's outline so every arrow lands on
%     the box edge rather than in empty space.
\draw[arr] (cat) -- (interp);
\draw[arr] (not) -- (interp);
\draw[arr] (bnd) -- (interp);

\draw[arr] (interp) -- (nec);
\draw[arr] (nec)    -- (dec);

% --- Burden -> interpretation (vertical, upward).
\draw[arr] (burden) -- (interp);

% --- Feedback loop:
%     (a) Privacy decision -> design implications  (vertical dashed on
%         the right).
%     (b) Design implications -> interpretive burden  (horizontal
%         dashed along the bottom).  With the narrower boxes there is
%         now ~3 cm of horizontal arrow visible between them.
\draw[fbarr] (dec)    -- (design);
\draw[fbarr] (design) -- (burden);

\end{tikzpicture}%
}
\caption{Conceptual model of the Necessity Gap. Three input streams
feed contextual interpretation, which converts them into a necessity
judgment and then a privacy decision (top row, solid arrows).
Interpretive burden is the cost that ambiguous categories and
boundaries impose on interpretation, and is the target of the
proposed design implications (bottom row). The dashed feedback loop
runs from the privacy decision down to the design implications, then
horizontally to the interpretive burden, which feeds back upward into
the interpretation.}
\label{fig:discussion-model}
\end{figure*}
